\subsection{Hilbert schemes and the exact positive theta input}

\paragraph{P-F53 (a geometric coefficient model for the identified input).}
Let $S$ be a complex projective K3 surface, and let $S^{[m]}$
be its Hilbert scheme of length-$m$ subschemes, with
$S^{[0]}=\operatorname{Spec}\mathbb C$.
Retain the positive-definite lattice $K=E_6$, its discriminant
cosets $\mu\in K^\vee/K$, and the quadratic form
$Q(x)=(x,x)/2$ of P-F46 and P-F51. Write the actual
Borcherds input from P-F46 as
\[
 g_\mu(\tau)=\frac{\Theta_{K+\mu}(\tau)}{\delta(\tau)}
        =\sum_a c_\mu(a)q^a,
 \qquad \delta(\tau)=q\prod_{r\ge1}(1-q^r)^{24}.
                                                        \tag{P-F53a}
\]
For each rational exponent $a$, define the following finite disjoint union:
\[
 X_{\mu,a}=
 \coprod_{\substack{x\in K+\mu\\
                  m=a+1-Q(x)\in\mathbb Z_{\ge0}}} S^{[m]}.
                                                        \tag{P-F53b}
\]
An empty disjoint union is allowed. Then $X_{\mu,a}$ is a
smooth projective complex scheme (possibly disconnected and
of different dimensions on its components), its odd rational
cohomology vanishes, and
\[
          c_\mu(a)=\dim_{\mathbb Q}H^*(X_{\mu,a},\mathbb Q).
                                                        \tag{P-F53c}
\]
Each $H^j(X_{\mu,a},\mathbb Q)$ is a pure Hodge structure of
weight $j$. The statement concerns purity separately in each
cohomological degree, and does not identify all the cohomology
with one pure weight or assert that it is Tate.

\paragraph{The cited Hilbert-scheme formula and the proof.}
The smoothness and projectivity of $S^{[m]}$ are the
surface Hilbert-scheme theorem recalled in G\"ottsche,
\emph{Hilbert Schemes of Points on Surfaces}, Section~1,
printed p.~484. His Betti-number formula is equation~(2.1),
printed p.~485. In his sign convention the Poincar\'e
polynomial is $p(Y,z)=\sum_j(-1)^jb_j(Y)z^j$.
For a K3 surface the Betti numbers are $(1,0,22,0,1)$,
and the formula specializes to
\[
 \sum_{m\ge0}p(S^{[m]},z)q^m
 =\prod_{r\ge1}
   \frac1{(1-z^{2r-2}q^r)
           (1-z^{2r}q^r)^{22}
           (1-z^{2r+2}q^r)}.
                                                        \tag{P-F53d}
\]
All exponents of $z$ on the right are even. Comparing
coefficients of $q^m z^{2j+1}$ gives
$b_{2j+1}(S^{[m]})=0$.
At $z=1$ the signed polynomial therefore equals the
total cohomology dimension, and
\[
 \sum_{m\ge0}\dim H^*(S^{[m]},\mathbb Q)q^m
                   =\prod_{r\ge1}(1-q^r)^{-24}.
                                                        \tag{P-F53e}
\]
The same Euler-number specialization is explicitly
Corollary~2.4, printed p.~62, of G\"ottsche,
\emph{Invariants of Moduli Spaces and Modular Forms};
the K3 value $e(S)=24$ is stated there on p.~66.

For any fixed $a$, a summand in (P-F53b) requires
$Q(x)\le a+1$. Positive definiteness of $K$ makes that
set of coset vectors finite. For each such vector the
integer $m$ is unique. Hence (P-F53b) is a finite
disjoint union of smooth projective schemes. Finite
disjoint unions preserve smoothness and projectivity;
their rational cohomology is the direct sum of the
component cohomologies. This proves odd-degree
vanishing and the asserted degreewise purity, using
the usual pure Hodge structure on smooth projective
cohomology.

Finally, multiplying (P-F53e) by
$q^{-1}\sum_{x\in K+\mu}q^{Q(x)}$ and taking the
coefficient of $q^a$ gives precisely the direct sum
in (P-F53b), so proves (P-F53c). All coefficient sums
are finite; there is no interchange of conditionally
convergent series. This completes the proof of the
coefficient model.\hfill$\square$

\paragraph{The completed lambda-ring identity.}
This construction also gives an exact, limited lambda-ring
dictionary for the inverse-discriminant operation. Let
$V=H^*(S,\mathbb Q)$ after forgetting its cohomological
grading, so $\dim V=24$. In the completed Grothendieck
ring of finite-dimensional rational vector spaces, put
$\sigma_t(V)=\sum_{j\ge0}[\operatorname{Sym}^jV]t^j$.
Then
\[
 \sum_{m\ge0}[H^*(S^{[m]},\mathbb Q)]q^m
   =\prod_{r\ge1}\sigma_{q^r}(V)
   =\operatorname{Sym}\!\left(V\otimes q\mathbb Q[q]\right)
                                                        \tag{P-F53f}
\]
as classes graded by the exponent of $q$. The last term
denotes the graded class of the symmetric algebra, with
one copy of $V$ in each positive degree. To prove the
identity, the generating function of the dimensions
of symmetric powers of a 24-dimensional vector space
is $(1-t)^{-24}$: a monomial basis counts weak
compositions into 24 parts. Its product over $r$
is (P-F53e). Dimension identifies the Grothendieck
ring of finite-dimensional rational vector spaces
with $\mathbb Z$, so equality of these dimensions
proves the stated class equality in every grade.
Every grade receives only finitely many contributions.
Thus no infinite-dimensional object is being inserted
as one element of the uncompleted ring.

Equation (P-F53f) uses the symmetric operations
$\sigma_t=\lambda_{-t}^{-1}$ in the sense explained in
the Grothendieck plethysm bridge. It is an equality
after forgetting the cohomological grading. This
elementary proof does not assert a canonical
cohomological-ring isomorphism or compatibility with
the nonstandard algebra's multiplication. The much
stronger geometric Heisenberg construction is not
needed for this statement.

\paragraph{The exact principal coefficients and what has been identified.}
The previously checked counts for the $E_6$ cosets
make the nonpositive part completely geometric:
\[
 X_{0,-1}=\operatorname{Spec}\mathbb C,\qquad
 X_{0,0}=S\ \coprod\ 72\operatorname{Spec}\mathbb C,
 \qquad
 X_{\pm,-1/3}=27\operatorname{Spec}\mathbb C.
                                                        \tag{P-F53g}
\]
Indeed the zero coset has one vector of norm zero
and 72 vectors with $Q=1$, while each other coset
has 27 vectors of minimal $Q=2/3$, as proved in
P-F46. Their total dimensions are $1,96,27,27$.
In particular the weight $96/2=48$ of the identified
Borcherds output can be read as half the total
cohomology dimension of $X_{0,0}$.

This supplies smooth projective coefficient spaces
for the \emph{entire particular input} of P-F46,
and expresses its inverse-discriminant factor by
completed symmetric operations on an actual
surface cohomology. Together with P-F51, both
factors now have specified geometric sources.
It does not turn the regularized multiplicative
Borcherds lift into the supplied iterated definite
unitary lift, nor give a cohomological map from
these spaces onto the cubic multiplicity spaces.
The coefficient-growth obstruction of P-F46 still
rules out identifying the whole input with all
$T^{(5)}$ grades by a positive affine reindexing.
The pure Hodge pieces here also need not be Tate:
already $S\subset X_{0,0}$ has its nonzero
$(2,0)$ part. This example therefore adds an
actual Hilbert-scheme and lambda-operation source
to the remaining operation-comparison problem,
without claiming that purity alone solves it.

Primary sources for the cited formulae:
\url{https://arxiv.org/pdf/math/0304302},
Section~1 p.~484 and equation~(2.1) p.~485;
\url{https://rendiconti.dmi.units.it/volumi/41/06.pdf},
Corollary~2.4 p.~62 and the K3 specialization p.~66.

\paragraph{The known product as an exterior operation on these actual spaces.}
There is also an exact completed character-level description of the
\emph{output} operation. The infinite product itself is the cited
Borcherds theorem, not a new automorphy or denominator theorem.
The following construction specifies its coefficient spaces and its
completion in terms of (P-F53b).

Write the lattice of P-F46 as
\[
 M=U\oplus L,\qquad
 L=U'\oplus E_8(-1)^2\oplus A_2(-1),
 \qquad \operatorname{sig}(L)=(1,19).
 \tag{P-F53h}
\]
Choose a basis $z,z'$ of the first hyperbolic plane with
$(z,z')=1$ and both vectors isotropic. This is a divisor-one
cusp. The discriminant groups of $L$ and $M$ are identified by
the direct sum. Choose a Weyl chamber $W$ for the input $g$,
and a rational positive vector $Y\in W$. Choose a positive integer
$b$ such that $bY\in L$ and put
\[
 h(\lambda)=b(\lambda,Y)\in\mathbb Z,
 \qquad
 \Gamma_Y=\{0\}\cup\{\lambda\in L^\vee:h(\lambda)>0\}.
 \tag{P-F53i}
\]
Only this auxiliary height is used for the completion; it is
not the coefficient degree in a cubic invariant ring.
In the category of $\Gamma_Y$-graded rational vector spaces
which are finite dimensional in every bounded height, define
\[
 \mathcal E_W=
 \bigoplus_{\substack{\lambda\in L^\vee\\h(\lambda)>0}}
 H^*\!\left(X_{[\lambda],\lambda^2/2},\mathbb Q\right)
                         \otimes e^\lambda,
 \qquad
 E_W=\operatorname{ch}\mathcal E_W
     =\sum_{h(\lambda)>0}c_{[\lambda]}(\lambda^2/2)e^\lambda.
 \tag{P-F53j}
\]
The cohomological grading is forgotten in this formula.
All its summands are actual spaces, not virtual differences.
Let $\widehat R_Y$ be the ring of integral formal sums supported
on $\Gamma_Y$ in the symbols $e^\lambda$, with
$e^\lambda e^\nu=e^{\lambda+\nu}$,
whose support is finite in every bounded positive height and
whose height-zero term is an integer multiple of $e^0$.
The completed exterior class is well defined, and satisfies
\[
 \begin{split}
 D_W&=\sum_{j\ge0}(-1)^j
                  \operatorname{ch}\bigl(\Lambda^j\mathcal E_W\bigr)\\
    &=\prod_{h(\lambda)>0}
          (1-e^\lambda)^{c_{[\lambda]}(\lambda^2/2)}
       \in 1+\widehat R_{Y,>0}.
 \end{split}
 \tag{P-F53k}
\]
In the convergence neighborhood of this cusp, substitute
$e^\lambda=e^{2\pi i(\lambda,Z)}$. Borcherds,
\emph{Automorphic forms with singularities on Grassmannians},
Theorem~13.3(5), arXiv version printed pp.~48--49, gives
\[
 \Psi_M(g)(Z)=C_W e^{2\pi i(\rho_W,Z)}D_W(Z),
                       \qquad C_W\ne0,
 \tag{P-F53l}
\]
where $\rho_W$ is its Weyl vector. Thus, after its scalar and
Weyl prefactor are removed, the identified multiplicative
theta lift is precisely the completed exterior Euler character
of the specified smooth-projective coefficient cohomologies.
Moreover $D_W^2=\lambda_{-1}(\mathcal E_W\oplus\mathcal E_W)$
at this character level. Restoring the prefactor, this is
$\Psi_M(g)^2$, whose image is the six-variable discriminant
in P-F55. The normalized local expression $D_W$ itself is
not asserted to be a global automorphic form or a homogeneous
element to which the invariant-ring map applies.

\emph{Proof of the completion and specialization.}
The coefficients of $g$ vanish below exponent $-1$ by
(P-F46a)--(P-F46b). Consequently a nonzero summand in
(P-F53j) has $\lambda^2\ge-2$.
For a fixed height bound $R$, decompose
$\lambda=tY+v$ with $v\perp Y$. Then
\[
 0<t\le \frac{R}{bY^2},\qquad
       -v^2\le 2+t^2Y^2.
 \tag{P-F53m}
\]
The form on $Y^\perp$ is negative definite. These inequalities
bound $\lambda$ in an ordinary Euclidean norm, so the discrete
lattice $L^\vee$ has only finitely many possible vectors.
Each cohomology space in the resulting finite list is finite
dimensional by (P-F53b). This proves the required bounded-height
finiteness.

Every nonzero summand has integral height at least one.
In height at most $R$, therefore, at most $R$ exterior factors
can occur, and they can only come from the finite-dimensional
truncation just proved. Tensor products, exterior operations,
and all coefficients in (P-F53k) are consequently finite
computations at each height. On that truncation the usual
exterior-power formula for a direct sum of one-weight vector
spaces gives $(1-e^\lambda)^{\dim H^*X}$ for its weight
$\lambda$ factor. This proves (P-F53k) by taking the compatible
truncations. This argument establishes the algebraic completion;
analytic convergence near the cusp is supplied by the cited
Borcherds theorem.

For this split hyperbolic-plane cusp, the integer $N$ in
Theorem~13.3(5) is one. Its finite cusp prefactor is empty,
and the restriction condition in its double product selects
the single discriminant class $[\lambda]\in L^\vee/L$.
Its phase $(\delta,z')$ is integral, since a representative
of that class can be taken entirely in $L^\vee$.
For every weight with nonzero exponent, the condition
$(\lambda,W)>0$ is equivalent to $h(\lambda)>0$:
negative-norm weights define the chamber walls, and a
nonzero weight of nonnegative norm has constant pairing
sign on a component of the positive cone. Thus that
theorem's product is exactly (P-F53k), proving (P-F53l).
Multiplicativity of $\lambda_{-1}$ under direct sums proves
the squaring assertion.\hfill$\square$

\paragraph{Recovering the input and the precise infinite-dimensional limit.}
The completed exterior expression determines the entire
root-graded character, not just a few principal coefficients.
For $r\ge1$ let $\psi^r$ fix integral coefficients and send
$e^\lambda$ to $e^{r\lambda}$. In the rationalized completion,
\[
 \log D_W=-\sum_{r\ge1}\frac{\psi^r(E_W)}r,
 \qquad
 E_W=-\sum_{r\ge1}\frac{\mu(r)}r\psi^r(\log D_W).
 \tag{P-F53n}
\]
Here $\mu$ is the ordinary arithmetic M\"obius function.
The first identity follows by expanding each
$\log(1-e^\lambda)$ in (P-F53k). Substituting it into
the second gives, at index $k$, the coefficient
$k^{-1}\sum_{r\mid k}\mu(r)$, which is one for $k=1$
and zero otherwise. These manipulations are legitimate
height by height: $\psi^r$ multiplies every positive
height by $r$, so only finitely many terms contribute.
This proves the inverse identity in the exact completion,
without an analytic logarithm or branch choice.

In fact $E_W$ contains every Fourier coefficient of $g$.
Choose a hyperbolic basis $e,f$ of $U'$ with
$(e,f)=1$ and $(e,Y),(f,Y)>0$. Such signs can be chosen
simultaneously: writing $Y=\alpha e+\beta f+Y_-$,
the inequality $2\alpha\beta+Y_-^2>0$ forces
$\alpha\beta>0$. For a class $\mu$ choose a representative
$\nu\in (E_8(-1)^2\oplus A_2(-1))^\vee$.
For any exponent $a\equiv\nu^2/2\pmod{\mathbb Z}$ and
any nonzero $w$ in this negative definite integral lattice,
put
\[
 \lambda_k=e+
       \left(a-\frac{(\nu+kw)^2}{2}\right)f+\nu+kw.
 \tag{P-F53o}
\]
Its $f$ coefficient is integral, its discriminant class
is $\mu$, and $\lambda_k^2/2=a$.
The height is a quadratic polynomial in $k$ with positive
leading coefficient $-b(w^2/2)(f,Y)$, so it is positive
for all sufficiently large integers $k$.
The coefficient of $e^{\lambda_k}$ in (P-F53j) is
therefore exactly $c_\mu(a)$. Coefficients outside the
allowed congruence class are already zero by the Weil type.
This proves the asserted full input recovery. It is a
formal character recovery, and gives no bound on the
complexity of computing the required coefficients.

The completion cannot be replaced by a finite actual
or virtual positively weighted vector space with the
same ordinary exterior character. Indeed choose a
nonzero integral isotropic vector $\ell\in U'$ with
$h(\ell)>0$. Every $r\ell$, $r\ge1$, has norm zero
and trivial discriminant class, so (P-F46b) gives
\[
       [e^{r\ell}]E_W=c_0(0)=96\qquad(r\ge1).
 \tag{P-F53p}
\]
Thus $E_W$ has infinite support. If a finite virtual
positive-height character $V$ satisfied
$\lambda_{-1}(V)=D_W$, the inverse formula (P-F53n),
which applies equally to integral virtual coefficients,
would give $V=E_W$, a contradiction. This excludes this
particular finite ordinary graded-space interpretation;
it does not exclude a determinant of a nontrivial
endomorphism, an infinite operator determinant, or the
automorphic determinant in the supplied \emph{Remark}.

The input spaces and the multiplicative operation are
now specified in one completed lambda-ring dictionary.
The argument supplies no Lie bracket, differential,
geometric convolution, nonstandard canonical basis,
or comparison with the totally definite iterated lift's
chosen vectors and periods. The relevant obstruction
gate is now comparison of these explicit completed
operations with those constructions and with the
actual multiplicity maps; it is not the absence of a
coefficient-space model or of a formula for its exterior
operation. The positive dimensions in $E_W$ also do not
make the alternating coefficients of $D_W$ a positive
Hilbert series. The known generalized Kac--Moody
denominator interpretation of the ambient $\Phi_{12}$
is not silently transferred to its $E_6$ quasi-pullback.
The product theorem used above is available at
\url{https://arxiv.org/pdf/alg-geom/9609022},
Theorem~13.3(5), printed pp.~48--49.
